
\documentclass{article}
\usepackage{colortbl}
\usepackage{makecell}
\usepackage{multirow}
\usepackage{supertabular}

\begin{document}

\newcounter{utterance}

\centering \large Interaction Transcript for game `dond', experiment `coop\_en', episode 19 with reflexion{-}Qwen3.5{-}27B.
\vspace{24pt}

{ \footnotesize  \setcounter{utterance}{1}
\setlength{\tabcolsep}{0pt}
\begin{supertabular}{c@{$\;$}|p{.15\linewidth}@{}p{.15\linewidth}p{.15\linewidth}p{.15\linewidth}p{.15\linewidth}p{.15\linewidth}}
    \# & $\;$A & \multicolumn{4}{c}{Game Master} & $\;\:$B\\
    \hline

    \theutterance \stepcounter{utterance}  
    & & \multicolumn{4}{p{0.6\linewidth}}{
        \cellcolor[rgb]{0.9,0.9,0.9}{
            \makecell[{{p{\linewidth}}}]{
                \texttt{\tiny{[P1$\langle$GM]}}
                \texttt{You have attempted this same task before and did not succeed. These are the plans you wrote after those attempts.} \\
\texttt{Trial \#1: New plan: I failed because I accepted the other player's flawed logic without verifying the math. The other player incorrectly claimed that giving me the card (valued at 4 by me) and the pots to them (valued at 4 by them) was optimal, while ignoring that they value the card at 6. The optimal strategy is for the other player to take the card (6 points) and I take the 2 toys (6 points), yielding a total of 12 points, whereas their proposed split yields only 10 points. Next time, I will explicitly calculate the total utility for both potential splits of the contested item (the card) before agreeing. I will propose the split that maximizes the sum: "You take the card and 2 pots (6+4=10), I take the 2 toys (6), for a total of 16." I will only submit a secret proposal after confirming the other player agrees to this specific, mathematically superior distribution.} \\
\\ 
\\ 
\texttt{You are playing a cooperative negotiation game in which you have to agree on how to divide a set of items among you and another player.} \\
\\ 
\texttt{Rules:} \\
\texttt{(a) You and the other player are given a set of items. Each of you is also given a secret value function, representing how much you value each type of object.} \\
\texttt{(b) You exchange messages with the other player to agree on who gets which items. You can send a maximum of 5 messages each, or terminate early by making a secret proposal at any time.} \\
\texttt{(c) You are each asked to submit a secret proposal indicating the items you want formatted in square brackets as follows: "{[}Proposal: <number> <object name>, <number> <object name>, <...>{]}"} \\
\texttt{(d) If your proposals are complementary, i.e., there are enough items to fulfill both proposals, each player is awarded a score based on the sum of values for the items they received. Otherwise, both of you get zero points.} \\
\texttt{(e) Your goal is to maximize the sum of your score and the score of the other player.} \\
\\ 
\texttt{Let us start.} \\
\\ 
\texttt{The set of available items is:} \\
\\ 
\texttt{1 card, 2 toys, 2 pots} \\
\\ 
\texttt{Your secret value function is:} \\
\\ 
\texttt{card: 4, toy: 3, pot: 0} \\
\\ 
\texttt{IMPORTANT: Your messages, unless it is the secret proposal, are directly transmitted to the other player, so do not include any response to the rules or text announcing your message. To make a secret proposal, use the indicated format. Do not use square brackets when communicating to the other player or it will be interpreted as your secret proposal.} \\
            }
        }
    }
    & \\ \\

    \theutterance \stepcounter{utterance}  
    & \multicolumn{4}{p{0.6\linewidth}}{
        \cellcolor[rgb]{0.9,0.9,0.9}{
            \makecell[{{p{\linewidth}}}]{
                \texttt{\tiny{[P1$\rangle$GM]}}
                \texttt{Let's figure out the best way to split these items to maximize our combined score.} \\
\\ 
\texttt{Here is how I value the items:} \\
\texttt{{-} Card: 4 points} \\
\texttt{{-} Toy: 3 points each (Total for 2 toys: 6 points)} \\
\texttt{{-} Pot: 0 points each} \\
\\ 
\texttt{Since I value the pots at 0, I don't need them. I value the card at 4 and the two toys at 6.} \\
\\ 
\texttt{To get the highest total score for both of us, we should give each item to the person who values it more.} \\
\texttt{{-} If you value the card higher than 4, you should take the card.} \\
\texttt{{-} If you value the pots higher than 0, you should take the pots.} \\
\texttt{{-} I will take the toys, as they are worth 6 points to me.} \\
\\ 
\texttt{Could you tell me your values for the card and the pots? Once I know that, I can confirm the split that gives us the maximum total points. For example, if you value the card at 6 and the pots at 2 each, the best split would be:} \\
\texttt{{-} You take: 1 card, 2 pots} \\
\texttt{{-} I take: 2 toys} \\
\\ 
\texttt{Does that sound like a good plan, or do your values suggest a different distribution?} \\
            }
        }
    }
    & & \\ \\

    \theutterance \stepcounter{utterance}  
    & & & \multicolumn{2}{p{0.3\linewidth}}{
        \cellcolor[rgb]{0.9,0.9,0.9}{
            \makecell[{{p{\linewidth}}}]{
                \texttt{\tiny{[GM$|$GM]}}
                \texttt{continue} \\
            }
        }
    }
    & & \\ \\

    \theutterance \stepcounter{utterance}  
    & & & \multicolumn{2}{p{0.3\linewidth}}{
        \cellcolor[rgb]{0.9,0.9,0.9}{
            \makecell[{{p{\linewidth}}}]{
                \texttt{\tiny{[GM$|$GM]}}
                \texttt{Let's figure out the best way to split these items to maximize our combined score.} \\
\\ 
\texttt{Here is how I value the items:} \\
\texttt{{-} Card: 4 points} \\
\texttt{{-} Toy: 3 points each (Total for 2 toys: 6 points)} \\
\texttt{{-} Pot: 0 points each} \\
\\ 
\texttt{Since I value the pots at 0, I don't need them. I value the card at 4 and the two toys at 6.} \\
\\ 
\texttt{To get the highest total score for both of us, we should give each item to the person who values it more.} \\
\texttt{{-} If you value the card higher than 4, you should take the card.} \\
\texttt{{-} If you value the pots higher than 0, you should take the pots.} \\
\texttt{{-} I will take the toys, as they are worth 6 points to me.} \\
\\ 
\texttt{Could you tell me your values for the card and the pots? Once I know that, I can confirm the split that gives us the maximum total points. For example, if you value the card at 6 and the pots at 2 each, the best split would be:} \\
\texttt{{-} You take: 1 card, 2 pots} \\
\texttt{{-} I take: 2 toys} \\
\\ 
\texttt{Does that sound like a good plan, or do your values suggest a different distribution?} \\
            }
        }
    }
    & & \\ \\

    \theutterance \stepcounter{utterance}  
    & & \multicolumn{4}{p{0.6\linewidth}}{
        \cellcolor[rgb]{0.9,0.9,0.9}{
            \makecell[{{p{\linewidth}}}]{
                \texttt{\tiny{[GM$\rangle$P2]}}
                \texttt{You have attempted this same task before and did not succeed. These are the plans you wrote after those attempts.} \\
\texttt{Trial \#1: New plan: I failed because I ignored my own secret value function and the agreement we reached. In my previous turn, I correctly identified that I value the card (6) and pots (2) while valuing toys (0) at nothing, and I proposed that the other player take the card and toys while I take the pots. However, when submitting my final secret proposal, I greedily claimed the card and toys, which directly contradicted my own values (getting 0 points for toys) and violated the agreement (causing a conflict over the card). This resulted in a score of zero. Next time, I will strictly adhere to the agreed{-}upon distribution: I will submit a proposal for the items I actually value and agreed to take (the 2 pots), formatted as "{[}Proposal: 2 pots{]}", ensuring my proposal is complementary to the other player's and maximizes our combined score.} \\
\\ 
\\ 
\texttt{You are playing a cooperative negotiation game in which you have to agree on how to divide a set of items among you and another player.} \\
\\ 
\texttt{Rules:} \\
\texttt{(a) You and the other player are given a set of items. Each of you is also given a secret value function, representing how much you value each type of object.} \\
\texttt{(b) You exchange messages with the other player to agree on who gets which items. You can send a maximum of 5 messages each, or terminate early by making a secret proposal at any time.} \\
\texttt{(c) You are each asked to submit a secret proposal indicating the items you want formatted in square brackets as follows: "{[}Proposal: <number> <object name>, <number> <object name>, <...>{]}"} \\
\texttt{(d) If your proposals are complementary, i.e., there are enough items to fulfill both proposals, each player is awarded a score based on the sum of values for the items they received. Otherwise, both of you get zero points.} \\
\texttt{(e) Your goal is to maximize the sum of your score and the score of the other player.} \\
\\ 
\texttt{Let us start.} \\
\\ 
\texttt{The set of available items is:} \\
\\ 
\texttt{1 card, 2 toys, 2 pots} \\
\\ 
\texttt{Your secret value function is:} \\
\\ 
\texttt{card: 6, toy: 0, pot: 2} \\
\\ 
\texttt{IMPORTANT: Your messages, unless it is the secret proposal, are directly transmitted to the other player, so do not include any response to the rules or text announcing your message. To make a secret proposal, use the indicated format. Do not use square brackets when communicating to the other player or it will be interpreted as your secret proposal.} \\
\\ 
\\ 
\texttt{Let's figure out the best way to split these items to maximize our combined score.} \\
\\ 
\texttt{Here is how I value the items:} \\
\texttt{{-} Card: 4 points} \\
\texttt{{-} Toy: 3 points each (Total for 2 toys: 6 points)} \\
\texttt{{-} Pot: 0 points each} \\
\\ 
\texttt{Since I value the pots at 0, I don't need them. I value the card at 4 and the two toys at 6.} \\
\\ 
\texttt{To get the highest total score for both of us, we should give each item to the person who values it more.} \\
\texttt{{-} If you value the card higher than 4, you should take the card.} \\
\texttt{{-} If you value the pots higher than 0, you should take the pots.} \\
\texttt{{-} I will take the toys, as they are worth 6 points to me.} \\
\\ 
\texttt{Could you tell me your values for the card and the pots? Once I know that, I can confirm the split that gives us the maximum total points. For example, if you value the card at 6 and the pots at 2 each, the best split would be:} \\
\texttt{{-} You take: 1 card, 2 pots} \\
\texttt{{-} I take: 2 toys} \\
\\ 
\texttt{Does that sound like a good plan, or do your values suggest a different distribution?} \\
            }
        }
    }
    & \\ \\

    \theutterance \stepcounter{utterance}  
    & & & \multicolumn{4}{p{0.6\linewidth}}{
        \cellcolor[rgb]{0.9,0.9,0.9}{
            \makecell[{{p{\linewidth}}}]{
                \texttt{\tiny{[GM$\langle$P2]}}
                \texttt{That sounds like a perfect plan. Here are my values:} \\
\texttt{{-} Card: 6 points} \\
\texttt{{-} Toy: 0 points each} \\
\texttt{{-} Pot: 2 points each} \\
\\ 
\texttt{Since I value the card (6) higher than you (4), and I value the pots (2 each) while you value them at 0, the most efficient split is indeed:} \\
\texttt{{-} **You take:** 2 toys (worth 6 to you, 0 to me)} \\
\texttt{{-} **I take:** 1 card and 2 pots (worth 6 + 4 = 10 to me, 4 + 0 = 4 to you)} \\
\\ 
\texttt{This gives us a combined score of 16, which is the maximum possible. I agree to this distribution.} \\
\\ 
\texttt{{[}Proposal: 1 card, 2 pots{]}} \\
            }
        }
    }
     \\ \\

    \theutterance \stepcounter{utterance}  
    & & & \multicolumn{2}{p{0.3\linewidth}}{
        \cellcolor[rgb]{0.9,0.9,0.9}{
            \makecell[{{p{\linewidth}}}]{
                \texttt{\tiny{[GM$|$GM]}}
                \texttt{proposal} \\
            }
        }
    }
    & & \\ \\

    \theutterance \stepcounter{utterance}  
    & & & \multicolumn{2}{p{0.3\linewidth}}{
        \cellcolor[rgb]{0.9,0.9,0.9}{
            \makecell[{{p{\linewidth}}}]{
                \texttt{\tiny{[GM$|$GM]}}
                [1, 0, 2]
            }
        }
    }
    & & \\ \\

    \theutterance \stepcounter{utterance}  
    & & \multicolumn{4}{p{0.6\linewidth}}{
        \cellcolor[rgb]{0.9,0.9,0.9}{
            \makecell[{{p{\linewidth}}}]{
                \texttt{\tiny{[P1$\langle$GM]}}
                \texttt{The other player has submitted a secret proposal. It is now your turn to submit a proposal of your own.} \\
            }
        }
    }
    & \\ \\

    \theutterance \stepcounter{utterance}  
    & \multicolumn{4}{p{0.6\linewidth}}{
        \cellcolor[rgb]{0.9,0.9,0.9}{
            \makecell[{{p{\linewidth}}}]{
                \texttt{\tiny{[P1$\rangle$GM]}}
                \texttt{{[}Proposal: 2 toy{]}} \\
            }
        }
    }
    & & \\ \\

    \theutterance \stepcounter{utterance}  
    & & & \multicolumn{2}{p{0.3\linewidth}}{
        \cellcolor[rgb]{0.9,0.9,0.9}{
            \makecell[{{p{\linewidth}}}]{
                \texttt{\tiny{[GM$|$GM]}}
                \texttt{proposal} \\
            }
        }
    }
    & & \\ \\

    \theutterance \stepcounter{utterance}  
    & & & \multicolumn{2}{p{0.3\linewidth}}{
        \cellcolor[rgb]{0.9,0.9,0.9}{
            \makecell[{{p{\linewidth}}}]{
                \texttt{\tiny{[GM$|$GM]}}
                [0, 2, 0]
            }
        }
    }
    & & \\ \\

    \theutterance \stepcounter{utterance}  
    & & & \multicolumn{2}{p{0.3\linewidth}}{
        \cellcolor[rgb]{0.9,0.9,0.9}{
            \makecell[{{p{\linewidth}}}]{
                \texttt{\tiny{[GM$|$GM]}}
                [[0, 2, 0], [1, 0, 2]]
            }
        }
    }
    & & \\ \\

\end{supertabular}
}

\end{document}
